\section{Go-To eksempler til ligningsløsning}

Her vil vi vise nogle generelle eksempler, man kan følge til mange almindelige ligninger.

\subsection*{Ligning 1}
Et simpelt udtryk:

\[
x + 5 = 2
\]

Vi ønsker at isolere $x$. Vi trækker derfor 5 fra på begge sider:

\[
x + 5 - 5 = 2 - 5 \Rightarrow x = -3
\]

---

\subsection*{Ligning 2}
Et udtryk med et scaleret $x$:

\[
5x - 2 = 1
\]

Vi starter med at fjerne alle konstanter på den side, $x$ står:

\[
5x - 2 + 2 = 1 + 2 \Rightarrow 5x = 3
\]

Omskrevet som brøkform:

\[
\frac{5x}{1} = \frac{3}{1}
\]

For at isolere $x$ dividerer vi med 5, dvs. ganger med 5 i nævneren på begge sider:

\[
\frac{5x}{5 \cdot 1} = \frac{3}{5 \cdot 1} \Rightarrow \frac{x}{1} = \frac{3}{5}
\]

Så vi har:

\[
x = \frac{3}{5}
\]

---

\subsection*{Ligning 3}
Et eksempel hvor $x$ står i nævneren:

\[
\frac{3}{x} + 2 = 1
\]

Først flytter vi konstanter over:

\[
\frac{3}{x} + 2 - 2 = 1 - 2 \Rightarrow \frac{3}{x} = -1
\]

Vi omskriver til brøkform:

\[
\frac{3}{x} = \frac{-1}{1}
\]

Vi vender brøken:

\[
\frac{x}{3} = \frac{1}{-1}
\]

Fjerner 3 i nævneren ved at gange med 3 på begge sider:

\[
\frac{3 \cdot x}{3} = \frac{3 \cdot 1}{-1} \Rightarrow \frac{x}{1} = \frac{3}{-1}
\]

Anvender reglen \(\frac{-a}{1} = -a\):

\[
x = -3
\]

---

\subsection*{Ligning 4}
Et eksempel med $x$ på begge sider:

\[
3x + 1 = x
\]

Vi flytter $x$ fra højre til venstre:

\[
3x + 1 - x = x - x \Rightarrow 2x + 1 = 0
\]

Trækker 1 fra på begge sider:

\[
2x + 1 - 1 = 0 - 1 \Rightarrow 2x = -1
\]

Omskrives til brøkform:

\[
\frac{2x}{1} = \frac{-1}{1}
\]

Dividerer med 2 (ganger med 2 i nævneren):

\[
\frac{2x}{1 \cdot 2} = \frac{-1}{1 \cdot 2} \Rightarrow \frac{x}{1} = \frac{-1}{2}
\]

Vi har dermed:

\[
x = -\frac{1}{2}
\]